\documentclass[10pt]{article}
\usepackage[T1]{fontenc}
\usepackage[utf8]{inputenc}
\usepackage[french]{babel}
\usepackage{amsmath,amssymb,mathrsfs,mathpazo}
\usepackage{geometry}
\geometry{a4paper,margin=2cm}
\usepackage{graphicx,array,tabularx,tabularray,colortbl,multirow,multicol,booktabs,enumitem,pifont,ulem,fancybox,diagbox,cancel}
\usepackage{pstricks,pstricks-add,pst-plot,pst-tree,pst-node,pst-text,pst-func,pst-math,pst-bspline,pst-3dplot,pst-eucl,pst-solides3d}
\let\pscalculate\undefined
\usepackage{pst-bezier}
\usepackage{tikz,pgfplots,tkz-tab}
\usetikzlibrary{arrows,arrows.meta,patterns,calc,positioning,trees,decorations.markings,shapes.geometric,shapes.misc,shapes.arrows,intersections,decorations.pathmorphing,decorations.pathreplacing,backgrounds,shadows,babel}
\pgfplotsset{compat=1.18}
\usepgfplotslibrary{fillbetween}
\usepackage[np]{numprint}
\usepackage{listings,makecell,pas-tableur,twoopt}
\let\barre\undefined
\usepackage[pstricks]{tablvar}
\definecolor{darkgreen}{RGB}{0,100,0}
\providecommand{\red}{\color{red}}
\providecommand{\blue}{\color{blue}}
\providecommand{\green}{\color{green}}
\providecommand{\black}{\color{black}}
\providecommand{\white}{\color{white}}
\providecommand{\gray}{\color{gray}}
\providecommand{\cyan}{\color{cyan}}
\providecommand{\magenta}{\color{magenta}}
\providecommand{\yellow}{\color{yellow}}
\providecommand{\violet}{\color{violet}}
\providecommand{\orange}{\color{orange}}
\providecommand{\euro}{\texteuro}
\providecommand{\EUR}{\texteuro}
\providecommand{\ieme}{\textsuperscript{e}}
\providecommand{\degres}{\ensuremath{^\circ}}
\providecommand{\degree}{\ensuremath{^\circ}}
\providecommand{\vv}[1]{\overrightarrow{#1}}
\providecommand{\Coord}[1]{#1}
\providecommand{\decosix}{}
\providecommand{\starredbullet}{\ensuremath{\bullet}}
\pagestyle{empty}
\setlength{\parindent}{0pt}
\usepackage{adjustbox}\DeclareUnicodeCharacter{2236}{:}\DeclareUnicodeCharacter{2190}{\ensuremath{\leftarrow}}
\providecommand{\R}{\mathbb{R}}
\providecommand{\N}{\mathbb{N}}
\providecommand{\Z}{\mathbb{Z}}
\providecommand{\Q}{\mathbb{Q}}
\providecommand{\D}{\mathbb{D}}
\providecommand{\C}{\mathbb{C}}
\providecommand{\e}{{\mathrm{e}}}
\providecommand{\vect}[1]{\overrightarrow{#1}}
\providecommand{\vectt}[1]{\overrightarrow{\mathrm{#1}}}
\providecommand{\vvt}[1]{\overrightarrow{\mathrm{#1}}}
\providecommand{\np}[2][]{{#2~#1}}
\providecommand{\vv}[1]{{\overrightarrow{#1}}}
\renewcommand{\d}{{\mathrm{d}}}
\providecommand{\degres}{{^\circ}}
\providecommand{\degre}{{^\circ}}
\providecommand{\blue}{\color{blue}}
\providecommand{\red}{\color{red}}
\providecommand{\enskip}{\quad}
\providecommand{\llbracket}{\lbrack\!\lbrack}
\providecommand{\rrbracket}{\rbrack\!\rbrack}
\providecommand{\euro}{{\text{€}}}
\providecommand{\no}{{\text{n°}}}
\providecommand{\interval}[2]{{[#1;#2]}}
\providecommand{\intervFO}[2]{{[#1;#2[}}
\providecommand{\dotfill}{\ldots}

\providecommand{\ds}{\displaystyle}
\providecommand{\textding}[1]{#1}
\providecommand{\Oij}{\(\mathrm{O};\vec{\imath},\vec{\jmath}\)}
\providecommand{\Oijk}{\(\mathrm{O};\vec{\imath},\vec{\jmath},\vec{k}\)}
\providecommand{\Ouv}{\(\mathrm{O};\vec{u},\vec{v}\)}
\providecommand{\intervFF}[2]{\ensuremath{[#1;#2]}}
\renewcommand{\euro}{\ifmmode\text{\texteuro}\else\texteuro\fi}
\DeclareUnicodeCharacter{00B0}{\ifmmode^{\circ}\else\textdegree\fi}
\DeclareUnicodeCharacter{20AC}{\ifmmode\text{\texteuro}\else\texteuro\fi}
% Énoncés : Éducation nationale. Transcription et corrigés : APMEP, Denis Vergès et contributeurs. Découpage et classement : CoopMaths, Sébastien Lozano, Éric Elter et contributeurs. Adaptation de la présentation : Maths Explorer.
% CC BY-SA : https://coopmaths.fr/www/about/
% https://coopmaths.fr/alea/static/eam/2026/tex/eam_specifique_2026_04_polynesie_automatismes_cor.tex
\begin{document}
\noindent\textit{APMEP / CoopMaths. CC BY-SA. Présentation adaptée par Maths Explorer.}\par\medskip
\subsection*{Question 1 \quad Réponse D}

On applique l'identité remarquable $(a-b)^2 = a^2 - 2ab + b^2$ avec $a = 5x$ et $b = 2$ :
\[(5x-2)^2 = (5x)^2 - 2\times 5x\times 2 + 2^2 = 25x^2 - 20x + 4.\]

\subsection*{Question 2 \quad Réponse D}

La multiplication est prioritaire sur la soustraction :
\[\dfrac{1}{4}\times\dfrac{8}{5} - \dfrac{1}{5} = \dfrac{8}{20} - \dfrac{1}{5} = \dfrac{2}{5} - \dfrac{1}{5} = \dfrac{1}{5}.\]

\subsection*{Question 3 \quad Réponse A}

Multiplier par $\dfrac{3}{4} = 0,75$ revient à appliquer le coefficient multiplicateur
$0,75 = 1 - 0,25$ : c'est une baisse de 25\,\%.

\subsection*{Question 4 \quad Réponse B}

Un produit de facteurs est nul si et seulement si l'un au moins de ses facteurs est nul :
\[(2x-4)(-3x+3) = 0 \iff 2x - 4 = 0 \ \text{ ou } \ -3x + 3 = 0 \iff x = 2 \ \text{ ou } \ x = 1.\]
D'où $S = \{2 ; 1\}$.

\subsection*{Question 5 \quad Réponse C}

Une baisse de 50\,\% correspond au coefficient $0,5$ et une hausse de 50\,\% au
coefficient $1,5$. Le coefficient multiplicateur global est leur produit :
\[0,5 \times 1,5 = 0,75 = 1 - 0,25.\]
Le prix a donc diminué de 25\,\%.

\subsection*{Question 6 \quad Réponse B}

Le coefficient directeur de la droite (AB) vaut
\[a = \dfrac{y_B - y_A}{x_B - x_A} = \dfrac{3 - (-1)}{1 - 3} = \dfrac{4}{-2} = -2.\]
On a donc $f(x) = -2x + b$, et comme B$(1~;~3)$ appartient à la droite :
$3 = -2\times 1 + b$, soit $b = 5$. D'où $f(x) = -2x + 5$.

\subsection*{Question 7 \quad Réponse C}

La moyenne coefficientée est la somme des notes multipliées par leur coefficient, divisée par
la somme des coefficients :
\[\overline{x} = \dfrac{18\times 1 + 11\times 2 + 16\times 1}{1 + 2 + 1} = \dfrac{18 + 22 + 16}{4} = \dfrac{56}{4} = 14.\]

\subsection*{Question 8 \quad Réponse D}

6,4~h, c'est 6~h plus 0,4~h. Or une heure compte 60 minutes :
\[0,4 \times 60 = 24.\]
Donc 6,4~h $=$ 6~h 24~min.

\end{document}
